\chapter{Commutativity via anti-involutions}\label{chap:commutativity}

A Hecke pair may carry an order-reversing symmetry that fixes each of its double cosets; when
it does, the structure constants are symmetric and the Hecke ring is commutative. Applying this
to the transpose on the arithmetic $\mathrm{GL}_n$ pair shows that the $\mathrm{GL}_n$ Hecke
algebra is commutative.

\begin{definition}\label{def:anti-involution}
\lean{HeckeRing.AntiInvolution}\leanok
\uses{def:hecke-pair}
Let $(G,H,\Delta)$ be a Hecke pair. An \emph{anti-involution} of $(G,H,\Delta)$ is an
anti-homomorphism $g\mapsto\bar g$ of $G$ that is involutive and preserves the pair. Explicitly,
it satisfies
\[
  \bar{\bar g}=g,\qquad \overline{ab}=\bar b\,\bar a,
\]
and maps $H$ into $H$ and $\Delta$ into $\Delta$, that is, $\bar g\in H$ whenever $g\in H$ and
$\bar g\in\Delta$ whenever $g\in\Delta$ [Sh, \S3.1].
\end{definition}

\begin{definition}\label{def:onHeckeCoset}
\lean{HeckeRing.AntiInvolution.onHeckeCoset}\leanok
\uses{def:anti-involution, def:hecke-coset}
Let $g\mapsto\bar g$ be an anti-involution of the Hecke pair $(G,H,\Delta)$. It induces a map on
the set $H\backslash\Delta/H$ of double cosets,
\[
  HgH \;\longmapsto\; H\bar gH .
\]
This is well defined: since the anti-involution reverses products and preserves $H$, two
representatives of the same double coset have barred images lying in a single double coset, so
the assignment depends only on $HgH$. The induced map is again involutive [Sh, \S3.1].
\end{definition}

\begin{theorem}\label{thm:mul-comm-of-ai}
\lean{HeckeRing.AntiInvolution.mul_comm_of_antiInvolution}\leanok
\uses{def:onHeckeCoset, def:hecke-multiplicity, thm:hecke-ring}
Let $g\mapsto\bar g$ be an anti-involution of the Hecke pair $(G,H,\Delta)$ whose induced map on
double cosets fixes every double coset, that is, $H\bar gH=HgH$ for all $g\in\Delta$. Then the
Hecke ring $\mathbb{T}$ is commutative: for all $f,g\in\mathbb{T}$,
\[
  f\cdot g \;=\; g\cdot f .
\]
\end{theorem}

\begin{proof}\leanok
\uses{def:onHeckeCoset, def:hecke-multiplicity, thm:hecke-ring}
Since the product on $\mathbb{T}$ is $\mathbb{Z}$-bilinear, it suffices to prove the identity on
a pair of basis elements $[D_1]$ and $[D_2]$, where the product is governed by the structure
constants $\mu(D_1,D_2;d)$. The claim therefore reduces to the symmetry
\[
  \mu(D_1,D_2;d)\;=\;\mu(D_2,D_1;d)
\]
for every double coset $d$. To see this, apply the anti-involution to a left-coset
decomposition $D_1D_2=\bigsqcup_i Hg_i$. Because $g\mapsto\bar g$ reverses products, it carries
this decomposition to one of $D_2D_1$; because it fixes each double coset, the double coset of
each $g_i$ is preserved. Tracking the left coset of a fixed representative $d$ under this
correspondence gives a bijection between the pairs of representatives counted by
$\mu(D_1,D_2;d)$ and those counted by $\mu(D_2,D_1;d)$, so the two multiplicities agree. Hence
the products $[D_1]\cdot[D_2]$ and $[D_2]\cdot[D_1]$ coincide on every basis pair, and
$\mathbb{T}$ is commutative [Sh, \S3.1, Prop.~3.8].
\end{proof}

\begin{definition}\label{def:commring-of-ai}
\lean{HeckeRing.instCommRing_of_antiInvolution}\leanok
\uses{thm:mul-comm-of-ai, thm:hecke-ring}
Given an anti-involution of the Hecke pair $(G,H,\Delta)$ whose induced map fixes every double
coset, the resulting commutativity of multiplication upgrades the ring structure on $\mathbb{T}$
to a \emph{commutative ring} structure, obtained from the ring structure on $\mathbb{T}$ by
supplying the symmetry $f\cdot g=g\cdot f$ from Theorem~\ref{thm:mul-comm-of-ai}.
\end{definition}

\begin{definition}\label{def:gl-pair-ai}
\lean{HeckeRing.GLn.GL_pair_antiInvolution}\leanok
\uses{def:anti-involution, def:GL-pair}
For the arithmetic $\mathrm{GL}_n$ Hecke pair, the \emph{transpose} $g\mapsto{}^{t}g$ is an
anti-involution. It reverses products, since ${}^{t}(ab)={}^{t}b\,{}^{t}a$, and is involutive,
since ${}^{t}({}^{t}g)=g$. Moreover it preserves the pair: transposing a matrix in
$\mathrm{SL}_n(\mathbb{Z})$ stays in $\mathrm{SL}_n(\mathbb{Z})$, and transposing a
positive-determinant integral matrix preserves both integrality of the entries and the value of
the determinant, so the positive-determinant integer monoid is preserved as well [Sh, \S3.1].
\end{definition}

\begin{lemma}\label{lem:gl-pair-fix}
\lean{HeckeRing.GLn.GL_pair_onHeckeCoset_eq}\leanok
\uses{def:gl-pair-ai, def:onHeckeCoset}
The transpose fixes every double coset of the arithmetic $\mathrm{GL}_n$ pair: for all $g$,
\[
  H\,{}^{t}gH \;=\; HgH .
\]
\end{lemma}

\begin{proof}\leanok
\uses{def:gl-pair-ai, def:onHeckeCoset}
Every double coset has a diagonal representative, given by the elementary divisors of $g$ via
the Smith normal form. A diagonal matrix is equal to its own transpose, so the transpose fixes
this representative and therefore fixes the whole double coset. Equivalently, a matrix and its
transpose share the same elementary divisors, hence lie in the same double coset [Sh, \S3.1].
\end{proof}

\begin{theorem}\label{thm:commring-hecke-algebra}
\lean{HeckeRing.GLn.instCommRing_HeckeAlgebra}\leanok
\uses{def:commring-of-ai, def:gl-pair-ai, lem:gl-pair-fix, def:hecke-algebra}
The $\mathrm{GL}_n$ Hecke algebra $\mathcal{H}_n$ is commutative.
\end{theorem}

\begin{proof}\leanok
\uses{def:commring-of-ai, def:gl-pair-ai, lem:gl-pair-fix, def:hecke-algebra}
The transpose is an anti-involution of the arithmetic $\mathrm{GL}_n$ pair
(Definition~\ref{def:gl-pair-ai}) whose induced map fixes every double coset
(Lemma~\ref{lem:gl-pair-fix}). Applying the commutativity criterion of
Definition~\ref{def:commring-of-ai} to this anti-involution endows $\mathcal{H}_n$ with a
commutative ring structure [Sh, \S3.1, Prop.~3.8].
\end{proof}
